Considérons la formule de quadrature suivante :
$$
\int_{-1}^1 f(x)\, \mathrm{d}x \simeq \omega_1 f(-\alpha) + \omega_2 f(\alpha)
$$
où $\alpha \in\, ]0\, ;\, 1]$.

\begin{enumerate}
    \item Déterminer les poids $\omega_1$ et $\omega_2$ pour que cette formule de quadrature soit exacte pour les polynômes de $\mathbb{R}_1[X]$.
    
    \item On adopte désormais les poids déterminés à la question précédente. Quelle est la formule obtenue lorsque $\alpha = 1$ ? Est-elle exacte pour les polynômes de $\mathbb{R}_2[X]$ ?
    
    \item Montrer que cette formule de quadrature est exacte sur $\mathbb{R}_2[X]$ pour une et une seule valeur de $\alpha$, à déterminer.
    
    \item On adopte la valeur de $\alpha$ déterminée à la question précédente. La formule est-elle exacte sur $\mathbb{R}_3[X]$ ? sur $\mathbb{R}_4[X]$ ?
    
    \item Adapter la formule obtenue à une intégrale $\int_0^1 f(x)\, \mathrm{d}x$, puis à une intégrale $\int_a^b f(x)\, \mathrm{d}x$.
    
    \item Exprimer la formule composite obtenue lorsque l'on subdivise l'intervalle $[a, b]$ en $N$ sous-intervalles.
\end{enumerate}